\documentclass{article}
\usepackage{graphicx}
\usepackage[margin=1.15in]{geometry}
\usepackage{xcolor}
\usepackage{parskip}
\usepackage{float}
\usepackage{fancyhdr}
\usepackage{array}
\usepackage{booktabs}
\usepackage{enumitem}
\usepackage{amsmath}

\pagestyle{fancy}
\fancyhf{}
\cfoot{\thepage}
\rhead{Lecture 20}
\lhead{MA 214: Applied Statistics (Spring 2026)}
\renewcommand{\headrulewidth}{.4mm}

\newcommand{\blankline}[1]{\underline{\hspace{#1}}}

\usepackage{listings}
\lstset{
  basicstyle=\ttfamily\small,
  columns=fullflexible,
  frame=single,
  framerule=0.4pt,
  breaklines=true,
  showstringspaces=false,
  keywordstyle=\color{blue!60!black},
  commentstyle=\color{green!40!black},
  stringstyle=\color{orange!60!black}
}

\begin{document}

\noindent\textbf{Name:}\blankline{2.25in}\hfill\textbf{BUID:}\blankline{1.55in}

\medskip

\noindent\textbf{Name:}\blankline{2.25in}\hfill\textbf{BUID:}\blankline{1.55in}\newline

\begin{center}
 \Large In-Class Worksheet: Posterior Inference \& Prediction
\end{center}

\subsection*{Scenario}
A field station on a coastal reserve is monitoring a rare shorebird during spring migration. From past years the morning count averaged about $2$ birds, so the team encodes that baseline as a prior on the daily rate $\lambda$:
\[
\lambda \sim \text{Gamma}(4,\,2) \qquad (\text{prior mean } 4/2 = 2 \text{ birds/morning}).
\]
This spring they record counts on $n = 6$ mornings: $3,\ 5,\ 2,\ 4,\ 6,\ 4$ (total $24$). The Gamma prior is conjugate to the Poisson count model, so the posterior is closed-form:
\[
\lambda \mid \text{data} \sim \text{Gamma}\!\left(4 + 24,\; 2 + 6\right) = \text{Gamma}(28,\, 8), \qquad \text{posterior mean } 28/8 = 3.5.
\]

\medskip
\noindent\textbf{Goals:}
\textbf{(1)} build a credible interval and test a hypothesis from this posterior,
\textbf{(2)} predict a new morning's count and check the model,
\textbf{(3)} debug a posterior-prediction script.

\medskip
\noindent\textbf{Tools.} You only need a few R one-liners: \texttt{qgamma}, \texttt{pgamma}, \texttt{rgamma}, \texttt{rpois}, and \texttt{quantile}. Work in \textbf{pairs}: one runs R, the other records answers.

\subsection*{Part 1: Credible Interval and Hypothesis Test}
\begin{enumerate}[label={\bfseries (\Alph*)},itemsep=.6em]
    \item \textbf{Credible interval.} The posterior is $\text{Gamma}(28, 8)$. Use \texttt{qgamma(c(0.025, 0.975), 28, 8)} to find the equal-tailed $95\%$ credible interval for $\lambda$. Write it down and complete the sentence: ``Given these data, there is a $95\%$ probability that the true morning rate $\lambda$ is between \rule{1.5cm}{0.4pt} and \rule{1.5cm}{0.4pt} birds.'' \\[3.5em]

    \item \textbf{Why not a confidence interval?} A classmate says ``there's a $95\%$ chance $\lambda$ is in that interval'' about a frequentist $95\%$ confidence interval. Why is that statement wrong for a confidence interval but correct for your credible interval? \\[4em]

    \newpage
    \item \textbf{Hypothesis test.} The team wants to know whether this spring's rate really beats the historical baseline of $2$ birds/morning: $H_0\!: \lambda \le 2$ vs.\ $H_1\!: \lambda > 2$. Compute $p_{\text{Bayes}} = P(\lambda \le 2 \mid \text{data})$ with \texttt{pgamma(2, 28, 8)}, then $s_{\text{Bayes}} = -\log_2 p_{\text{Bayes}}$.

    \vspace{0.3em}
    \begin{center}
    \renewcommand{\arraystretch}{1.4}
    \begin{tabular}{>{\bfseries}p{5cm}c}
    \toprule
    $p_{\text{Bayes}} = P(\lambda \le 2 \mid \text{data})$ & \blankline{3.5cm} \\
    $s_{\text{Bayes}} = -\log_2 p_{\text{Bayes}}$ (bits) & \blankline{3.5cm} \\
    \bottomrule
    \end{tabular}
    \end{center}
    \vspace{0.4em}

    Is there strong evidence the rate is up this spring? What does the number of bits mean in plain language? \\[3.5em]
\end{enumerate}

\subsection*{Part 2: Prediction and Predictive Checking}
\begin{enumerate}[label={\bfseries (\Alph*)},itemsep=.6em]
    \item \textbf{Posterior prediction.} You want to predict tomorrow morning's count $\tilde{y}$. In R, draw $M = 10{,}000$ posterior--predictive values:
    \begin{center}
    \texttt{lam <- rgamma(10000, 28, 8);\quad y <- rpois(10000, lam)}
    \end{center}
    Report the predictive mean and the $95\%$ predictive interval \texttt{quantile(y, c(0.025, 0.975))}. \\[3em]

    \item \textbf{Averaging vs.\ plug-in.} A colleague instead predicts with a single number: \texttt{rpois(10000, 3.5)} (plug in the posterior mean $\hat\lambda = 3.5$). Which prediction is wider, and \emph{why}? What source of uncertainty does the plug-in version leave out? \\[4em]

    \item \textbf{Predictive check.} To test the Poisson model, you simulate $5000$ replicate weeks of $6$ mornings each (draw $\lambda$ from the posterior, then $6$ Poisson counts), and record the \emph{busiest morning} in each replicate. The observed busiest morning was $6$ birds. The fraction of replicates with a busiest morning $\ge 6$ came out to $0.30$. Is there any reason to doubt the model? What predictive $p$-value \emph{would} have worried you? \\[4em]
\end{enumerate}

\newpage

\subsection*{Part 3: Find the Bugs}
A teammate wrote the script below to (i) report a $95\%$ credible interval for $\lambda$ and (ii) predict tomorrow's count by posterior prediction. It runs without error, but \textbf{three lines are statistically wrong}. The posterior is $\text{Gamma}(28, 8)$ and the prior was $\text{Gamma}(4, 2)$.

\begin{lstlisting}[language=R]
set.seed(1)
M          <- 100000
alpha_post <- 28    # posterior Gamma(28, 8)
beta_post  <- 8
alpha_prior <- 4    # prior    Gamma(4, 2)
beta_prior  <- 2

# --- 95% credible interval for lambda ---
lam_draws <- rgamma(M, alpha_prior, beta_prior)
ci <- quantile(lam_draws, c(0.05, 0.95))
cat("95% credible interval:", round(ci, 2), "\n")

# --- posterior prediction: tomorrow's count ---
lambda <- rgamma(1, alpha_post, beta_post)
y_pred <- rpois(M, lambda)
pi <- quantile(y_pred, c(0.025, 0.975))
cat("95% prediction interval:", pi, "\n")
\end{lstlisting}

\begin{enumerate}[label={\bfseries (\Alph*)},itemsep=.6em]
    \item \textbf{Bug 1 (credible interval, wrong distribution).} Which line draws from the \emph{wrong} distribution, and what should it be? What would the reported interval describe as written? \\[3.5em]

    \item \textbf{Bug 2 (credible interval, wrong quantiles).} The interval is supposed to be $95\%$. What is wrong with \texttt{c(0.05, 0.95)}, and what confidence level does it actually give? \\[3em]

    \item \textbf{Bug 3 (prediction, plug-in instead of averaging).} The prediction lines draw a \emph{single} \texttt{lambda} and reuse it for all $M$ Poisson draws. Why does this defeat the purpose of posterior prediction, and how would you fix it so each draw uses its own $\lambda$? (Connect your answer to Part 2B.) \\[4em]

    \item \textbf{Fix it.} Rewrite the three wrong lines correctly. \\[5em]
\end{enumerate}

\end{document}
